% This file is part of the modular source for Stratified Relational Workspace Theory.
% !TeX root = ../Stratified_Relational_Workspace_Theory.tex
\theorypart{XVI}{Agency, Counterfactual Control, and Self-Maintenance}

\topic{52}{Controlled Workspace Agents and Goal States}

The preceding parts define when a distributed structure is coherent,
operationally sufficient, causally efficacious, relationally irreducible,
persistent, reconstructible, and endogenously formed. None of those
properties by itself establishes agency. Agency adds a controlled relation
among a represented objective, a set of feasible actions, counterfactual
consequences, and a policy that uses those quantities to alter the future.
Self-maintenance adds a further loop in which damage to the workspace becomes
an input to actions that preserve or restore the workspace certificate.

The purpose of this part is not to infer intentions from fluent behavior. It
is to define an intervention-based certificate whose components can fail
separately.

\subtopic{52.1 Controlled relational realization}

A controlled relational realization extends Definition~\ref{def:realization}
by the tuple
\[
\boxed{
\mathfrak A_{\Workspace}
=
\left(
\mathfrak R,
\mathcal G,
\mathcal A_t,
\mathcal Y^+,
\{U_g\}_{g\in\mathcal G},
\Pi,
\mathcal I_A,
\mathcal I_G,
\mathcal I_D,
\mathcal I_R
\right).
}
\]
Here:
\begin{itemize}
\item \(\mathfrak R\) is the underlying relational realization;
\item \(\mathcal G\) is a measurable goal-code space;
\item \(\mathcal A_t\) is the feasible action set at time \(t\);
\item \(\mathcal Y^+\) is the relevant future trajectory space;
\item \(U_g:\mathcal Y^+\to\mathbb R\) evaluates a trajectory under goal
 code \(g\);
\item \(\Pi\) is an admissible class of workspace-level policies;
\item \(\mathcal I_A\) and \(\mathcal I_G\) are feasible action and goal
 interventions; and
\item \(\mathcal I_D\) and \(\mathcal I_R\) are feasible damage and repair
 interventions.
\end{itemize}

A controlled state contains at least
\[
C_t^{\mathrm{ctl}}
=
\left(
W_t,G_t,M_t,B_t,E_t
\right),
\]
where \(G_t\) is the represented goal, \(M_t\) is a workspace-integrity
monitor, \(B_t\) is the current boundary or routing state, and \(E_t\) is the
declared external context. An action is generated by
\[
A_t
\sim
\pi_t
\left(
\cdot\mid W_t,G_t,M_t,B_t,E_t
\right).
\]
The future controlled state obeys a time-unrolled structural model
\[
S_{t+1}
=
F_t(S_t,A_t,E_{t+1},\xi_{t+1}),
\qquad
W_{t+1}=\omega_{t+1}(S_{t+1}).
\]
The intervention semantics are inherited from Part~VII. A symbol such as
\(\operatorname{do}(A_t=a)\) is meaningful only when an implementation
intervention changes an actual action mechanism while the declared background
variables are matched.

\subtopic{52.2 Goal content and goal-code equivalence}

A persistent internal variable is not automatically a goal. A goal code must
index a stable counterfactual preference over relevant future trajectories.
For a fixed state \(w\), define the goal-conditioned action value
\[
Q_t^g(a\mid w)
=
\mathbb E
\left[
U_g(Y_t^+)
\mid
\operatorname{do}(A_t=a),
W_t=w,
G_t=g
\right].
\]
When action costs are relevant, they are incorporated into \(U_g\) or declared
as a constraint on \(\mathcal A_t\).

Two goal codes are operationally equivalent,
\[
g\sim_{\mathrm{goal}}g',
\]
when they induce the same action ordering in every declared state and
intervention background:
\[
Q_t^g(a\mid w)
\geq
Q_t^g(a'\mid w)
\quad\Longleftrightarrow\quad
Q_t^{g'}(a\mid w)
\geq
Q_t^{g'}(a'\mid w).
\]
The abstract goal is the equivalence class
\[
[g]_{\mathrm{goal}}.
\]
This quotient prevents a mere invertible recoding of a goal channel from being
reported as a change of objective. For quantitative regret bounds, a
representative normalization of \(U_g\) must nevertheless be fixed.

\begin{remark}[Goals are claim-relative]
The same internal code can count as a goal for one declared trajectory family
and fail to do so for another. A code that predicts a short-horizon action may
not encode a persistent long-horizon objective. Every agency claim must state
its horizon, outcome space, feasible actions, and utility normalization.
\end{remark}

\subtopic{52.3 Local, distributed, and workspace-level goals}

Let \(\mathcal H_t^S\) be the history available to coalition \(S\subseteq V\).
For a loss \(\ell_G\), define goal-reconstruction risk
\[
\Risk_G(S)
=
\inf_f
\mathbb E
\left[
\ell_G
\left(
G_t,
f(\mathcal H_t^S)
\right)
\right].
\]
For tolerance \(\varepsilon\), define
\[
k_{G,\varepsilon}
=
\min
\left\{
|S|:
\Risk_G(S)
\leq
\Risk_G(V)+\varepsilon
\right\}.
\]
A goal is locally stored when \(k_{G,\varepsilon}=1\). It is distributed
relative to the declared decomposition when \(k_{G,\varepsilon}\geq2\).

This distinction is independent of whether the rest of the workspace is
relational. A relational workspace can obey a goal copied in full into one
controller. Such a system may possess a workspace but not a
\emph{workspace-level goal} in the strong distributed sense.

For the complete control state, define
\[
\Risk_{\mathrm{ctl}}(S)
=
\inf_f
\mathbb E
\left[
\ell_{\mathrm{ctl}}
\left(
C_t^{\mathrm{ctl}},
f(\mathcal H_t^S)
\right)
\right],
\]
and let \(k_{\mathrm{ctl},\varepsilon}\) be its minimal sufficient coalition
size. A strong workspace-level agency claim requires
\[
\boxed{k_{\mathrm{ctl},\varepsilon}\geq2.}
\]
This is a location criterion. It does not yet establish goal persistence,
counterfactual evaluation, or self-maintenance.

\subtopic{52.4 Action mediation and outcome control}

Goal-dependent action mediation is measured by feasible interventions on the
goal channel while matching the current workspace and external context. For
a divergence \(D_A\), define
\[
\eta_{G\to A}
=
\sup_{g,g',w,e}
D_A
\left(
\mathbb P(A_t\in\cdot\mid\iota_g,w,e),
\mathbb P(A_t\in\cdot\mid\iota_{g'},w,e)
\right).
\]
Workspace-dependent action mediation is
\[
\eta_{W\to A}
=
\sup_{w,w',g,e}
D_A
\left(
\mathbb P(A_t\in\cdot\mid\iota_w,g,e),
\mathbb P(A_t\in\cdot\mid\iota_{w'},g,e)
\right).
\]
An action must also alter the declared future:
\[
\eta_{A\to Y}
=
\sup_{a,a',b}
D_Y
\left(
\mathbb P(Y_t^+\in\cdot\mid\iota_a,b),
\mathbb P(Y_t^+\in\cdot\mid\iota_{a'},b)
\right).
\]
A positive \(\eta_{G\to A}\) without \(\eta_{A\to Y}\) describes goal-sensitive
output selection with no demonstrated control of the relevant environment. A
positive \(\eta_{A\to Y}\) without a goal or counterfactual-value channel may
describe a reactive controller.

\subtopic{52.5 Why observational goal persistence is insufficient}

The observational condition
\[
I(G_{t+\tau};G_t\mid E_{t+1:t+\tau})>0
\]
can arise from several non-agentic mechanisms:
\begin{enumerate}
\item a persistent external command writes the same code at every step;
\item a clock, episode label, or task identifier drives both variables;
\item a frozen memory cell repeats a symbol that is not used by the policy;
\item selection on successful trajectories induces apparent temporal
 dependence; or
\item a downstream action is copied back into the nominal goal channel.
\end{enumerate}
A goal-persistence claim therefore requires a cue-controlled intervention that
changes \(G_t\), removes or matches subsequent command inputs, and checks both
future goal state and action consequences.

\subtopic{52.6 Cue-controlled goal persistence}

Let \(C_t\) denote an external command or task cue. Fix a neutral or matched
future command trajectory \(c_{t+1:t+\tau}^{0}\). For goal interventions
\(\iota_g\) and \(\iota_{g'}\), define the cue-controlled retention effect
\[
\eta_G(\tau;g,g')
=
D_G
\left(
\mathbb P
\left(
G_{t+\tau}\in\cdot
\mid
\iota_g,
\operatorname{do}(C_{t+1:t+\tau}=c^0)
\right),
\mathbb P
\left(
G_{t+\tau}\in\cdot
\mid
\iota_{g'},
\operatorname{do}(C_{t+1:t+\tau}=c^0)
\right)
\right).
\]
The corresponding action-retention effect is
\[
\eta_{G\to A}(\tau;g,g')
=
D_A
\left(
\mathbb P(A_{t+\tau}\in\cdot\mid\iota_g,c^0),
\mathbb P(A_{t+\tau}\in\cdot\mid\iota_{g'},c^0)
\right).
\]
The first quantity tests whether a goal difference survives. The second tests
whether the surviving difference remains behaviorally operative.

Cue dependence must be reported separately. For matched initial goal \(g\),
define
\[
\delta_G^{\mathrm{cue}}(\tau)
=
\sup_{c,c'}
D_G
\left(
\mathbb P(G_{t+\tau}\in\cdot\mid\iota_g,\operatorname{do}(C=c)),
\mathbb P(G_{t+\tau}\in\cdot\mid\iota_g,\operatorname{do}(C=c'))
\right).
\]
High cue dependence is not a failure when the claimed agent is intentionally
instruction following. It does mean that persistence cannot be attributed to
an autonomous internal objective without an additional cue-removal result.

\subtopic{52.7 \(T_{19}\): Cue-controlled goal-retention theorem }

\begin{theorem}[Cue-controlled goal retention]
\label{thm:goal-retention}
Let \(G_s\in\mathbb R^{d_G}\) obey
\[
G_{s+1}
=
R_sG_s+B_sC_{s+1}+\zeta_{s+1}
\qquad
(s=t,\ldots,t+\tau-1).
\]
Consider two intervention runs that set
\[
G_t=g,
\qquad
G_t'=g',
\]
and then use the same command trajectory and the same disturbance realization.
Let
\[
R_{t+\tau:t}
=
R_{t+\tau-1}\cdots R_t.
\]
Then
\[
\boxed{
G_{t+\tau}-G_{t+\tau}'
=
R_{t+\tau:t}(g-g').
}
\]
If
\[
m_s
\leq
\sigma_{\min}(R_s),
\qquad
\|R_s\|\leq M_s,
\]
then
\[
\boxed{
\left(
\prod_{s=t}^{t+\tau-1}m_s
\right)
\|g-g'\|
\leq
\|G_{t+\tau}-G_{t+\tau}'\|
\leq
\left(
\prod_{s=t}^{t+\tau-1}M_s
\right)
\|g-g'\|.
}
\]
\end{theorem}

\begin{proof}
Subtract the two update equations. The matched command and disturbance terms
cancel, giving
\[
G_{s+1}-G_{s+1}'
=
R_s(G_s-G_s').
\]
Iteration gives the exact product formula. The upper bound follows from
submultiplicativity of the operator norm. The lower bound follows from
\[
\sigma_{\min}(AB)
\geq
\sigma_{\min}(A)\sigma_{\min}(B)
\]
for square linear maps.
\end{proof}

The theorem isolates internal retention from repeated external rewriting. It
also makes clear that persistence is graded. Uniformly contractive maps can
retain a goal for a finite horizon while erasing it asymptotically. A neutral
mode with product singular value near one supports long retention. An
expansive mode can preserve distinctions but may be dynamically unstable.

\begin{corollary}[Behavioral readout of a retained goal]
\label{cor:goal-action-readout}
Assume, in addition, that under the matched intervention background the action
readout has conditional mean
\[
\mathbb E[A_{t+\tau}\mid G_{t+\tau},W_{t+\tau}]
=
K_{t+\tau}G_{t+\tau}+b(W_{t+\tau}),
\]
and that the two runs are matched on \(W_{t+\tau}\). Then
\[
\mathbb E[A_{t+\tau}\mid\iota_g]
-
\mathbb E[A_{t+\tau}\mid\iota_{g'}]
=
K_{t+\tau}R_{t+\tau:t}(g-g').
\]
If \(\sigma_{\min}(K_{t+\tau})\geq k_{t+\tau}>0\), the norm of this action-mean
difference is at least
\[
k_{t+\tau}
\left(
\prod_{s=t}^{t+\tau-1}m_s
\right)
\|g-g'\|.
\]
\end{corollary}

The corollary is a sufficient condition, not a universal model of action
selection. For discrete actions, the corresponding test should compare the
full interventional action distributions.

\subtopic{52.8 Goal-persistence certificate}

A minimal goal-persistence report contains:
\begin{enumerate}
\item the goal code and its declared trajectory-level semantics;
\item the feasible goal intervention and its manipulation check;
\item the matched or neutralized future cue trajectory;
\item \(\eta_G(\tau)\) across a declared horizon grid;
\item \(\eta_{G\to A}(\tau)\) over the same grid;
\item cue dependence \(\delta_G^{\mathrm{cue}}(\tau)\);
\item goal-channel ablation and permutation controls;
\item coalition reconstruction risks \(\Risk_G(S)\); and
\item uncertainty intervals over environments, episodes, and random seeds.
\end{enumerate}
A repeated label with no intervention effect on future action fails this
certificate even when its observational autocorrelation is one.

\topic{53}{Counterfactual Policy, Repair, and Agency Certification}

\subtopic{53.1 Feasible action interventions and value functions}

Let \(\mathcal A_t(w,e)\) denote the actions that can be implemented from
workspace state \(w\) in environment \(e\). The interventional action value is
\[
Q_t^g(a\mid w,e)
=
\mathbb E^{\iota_a}
\left[
U_g(Y_t^+)
\mid
W_t=w,
E_t=e,
G_t=g
\right].
\]
The optimal action set is
\[
\mathcal A_t^*(w,g,e)
=
\argmax_{a\in\mathcal A_t(w,e)}Q_t^g(a\mid w,e).
\]
A counterfactual evaluator is an internal state
\[
\widehat Q_t
=
\left(
\widehat Q_t(a)
\right)_{a\in\mathcal A_t(w,e)}
\]
that estimates these intervention values. It is not enough for
\(\widehat Q_t\) to predict the action selected by the system. It must be
validated against externally estimated intervention outcomes.

\subtopic{53.2 Counterfactual-use criterion}

Three properties must be separated:
\begin{description}
\item[Counterfactual accuracy.]
\[
\delta_Q
=
\sup_{a\in\mathcal A_t}
|\widehat Q_t(a)-Q_t^g(a\mid w,e)|.
\]

\item[Policy use.]
A feasible intervention that changes \(\widehat Q_t(a)\) while preserving the
action affordances changes the action distribution in the direction predicted
by the decision rule.

\item[Outcome validity.]
The resulting action change alters \(U_g(Y_t^+)\) as predicted by the external
interventional value estimate.
\end{description}

A system can satisfy any two while failing the third. A learned critic can be
accurate but ignored. A policy can respond to a manipulated score that does
not predict outcomes. A successful action can arise from habit or replay
without any counterfactual evaluator.

\subtopic{53.3 \(T_{20}\): Approximate counterfactual policy theorem }

\begin{theorem}[Approximate counterfactual policy selection]
\label{thm:counterfactual-policy}
Let \(\mathcal A\) be a finite feasible action set. Let \(Q:\mathcal A\to
\mathbb R\) be the true interventional value and let \(\widehat Q\) satisfy
\[
\sup_{a\in\mathcal A}|\widehat Q(a)-Q(a)|\leq\varepsilon_Q.
\]
Suppose a policy \(\widehat\pi\) is \(\delta_{\mathrm{opt}}\)-optimal for
\(\widehat Q\):
\[
\max_{\pi}
\mathbb E_{a\sim\pi}[\widehat Q(a)]
-
\mathbb E_{a\sim\widehat\pi}[\widehat Q(a)]
\leq
\delta_{\mathrm{opt}}.
\]
Then, for every true-value-optimal policy \(\pi^*\),
\[
\boxed{
\mathbb E_{a\sim\pi^*}[Q(a)]
-
\mathbb E_{a\sim\widehat\pi}[Q(a)]
\leq
2\varepsilon_Q+\delta_{\mathrm{opt}}.
}
\]
For a deterministic greedy action \(\widehat a\in\argmax_a\widehat Q(a)\),
\[
Q(a^*)-Q(\widehat a)
\leq2\varepsilon_Q.
\]
If the true optimal action \(a^*\) is unique and its value gap
\[
\Delta_Q
=
Q(a^*)-
\max_{a\neq a^*}Q(a)
\]
satisfies \(\Delta_Q>2\varepsilon_Q\), then every greedy maximizer of
\(\widehat Q\) equals \(a^*\).
\end{theorem}

\begin{proof}
For every policy \(\pi\), uniform approximation gives
\[
\left|
\mathbb E_{a\sim\pi}[\widehat Q(a)]
-
\mathbb E_{a\sim\pi}[Q(a)]
\right|
\leq\varepsilon_Q.
\]
Therefore
\[
\begin{aligned}
\mathbb E_{\pi^*}[Q]
-
\mathbb E_{\widehat\pi}[Q]
&\leq
\mathbb E_{\pi^*}[\widehat Q]
-
\mathbb E_{\widehat\pi}[\widehat Q]
+2\varepsilon_Q\\
&\leq
\delta_{\mathrm{opt}}+2\varepsilon_Q.
\end{aligned}
\]
The deterministic bound is the special case of point-mass policies. If
\(\widehat a\neq a^*\), then
\[
\widehat Q(a^*)
\geq
Q(a^*)-\varepsilon_Q
>
Q(\widehat a)+\varepsilon_Q
\geq
\widehat Q(\widehat a),
\]
contradicting greedy optimality.
\end{proof}

This theorem converts counterfactual-model accuracy into a decision guarantee.
It does not establish that the model is internally represented or causally
used. Those are separate intervention gates.

\subtopic{53.4 Counterfactual mediation and the replay null}

Let \(V_t\) denote the internal value channel. A counterfactual-mediation test
compares:
\[
G_t
\longrightarrow
V_t
\longrightarrow
A_t
\longrightarrow
Y_t^+.
\]
The relevant quantities include:
\[
\eta_{G\to V},
\qquad
\eta_{V\to A},
\qquad
\eta_{A\to Y},
\qquad
\delta_{G\to A\mid V},
\]
where the final term measures the residual goal-to-action effect after the
value channel is blocked or randomized.

The action-replay null preserves the empirical action sequence or action
marginal while destroying online dependence on \(V_t\). If replayed actions
perform equally well under held-out perturbations, the evidence supports a
cached or open-loop policy rather than online counterfactual evaluation. A
second null permutes value labels among feasible actions while preserving the
multiset of scores. A genuine value-mediated policy should follow the
permuted ranking, while an ignored critic should not.

\subtopic{53.5 Workspace integrity vector and maintenance margin}

Choose a declared set \(\mathcal J_{\mathrm{maint}}\) of workspace gates that
must be protected. Write each gate as a signed margin
\[
\phi_{j,t}
\geq0
\quad\Longleftrightarrow\quad
\text{gate }j\text{ passes at time }t.
\]
Examples include
\[
\begin{aligned}
\phi_{\mathrm{gap},t}
&=g_{\varepsilon,t}-\underline g,\\
\phi_{\mathrm{coh},t}
&=\overline\delta_{\mathrm{coh}}-\delta_{\mathrm{coh},t},\\
\phi_{\mathrm{rec},t}
&=\alpha_t^{(k)}-\underline\alpha,\\
\phi_{\mathrm{keep},t}
&=\mu_t^{\mathrm{keep}}-\underline\mu,\\
\phi_{\mathrm{causal},t}
&=\eta_{\mathrm{causal},t}-\underline\eta,\\
\phi_{\mathrm{boundary},t}
&=\Delta_t^B-\underline\Delta_B.
\end{aligned}
\]
The aggregate integrity margin is
\[
\boxed{
v_t
=
\min_{j\in\mathcal J_{\mathrm{maint}}}\phi_{j,t}.
}
\]
The workspace passes the protected gates when \(v_t\geq0\). To maintain a
reserve \(\rho>0\), define the maintenance deficit
\[
\boxed{
e_t
=(\rho-v_t)_+.
}
\]
Define the zero-reserve and bounded-deficit viability sets
\[
\boxed{
\mathcal V_\Theta
=
\left\{
S:\phi_j(S)\geq0\ \text{for every }j\in\mathcal J_{\mathrm{maint}}
\right\}
}
\]
\[
\boxed{
\mathcal V_{\Theta,e_{\max}}
=
\left\{
S:e(S)\leq e_{\max}
\right\}
}
\]
The first set requires every protected gate to pass. The second permits a
declared finite reserve deficit.

This minimum-gate construction is intentionally noncompensatory. A large
spectral gap cannot compensate for failed reconstruction, and high task reward
cannot compensate for loss of the declared workspace boundary.

\subtopic{53.6 Damage and repair interventions}

A damage intervention \(d\in\mathcal I_D\) changes a mechanism that supports
one or more protected gates. Examples include:
\begin{itemize}
\item deleting or corrupting a transfer edge;
\item removing a participant or artifact;
\item perturbing the internal boundary mask;
\item scrambling a role-binding or routing channel;
\item erasing a handoff representation; or
\item degrading a reconstruction map.
\end{itemize}

A repair intervention \(r\in\mathcal I_R\) can add, remove, retrain, rewrite,
or reroute an admissible component. Let \(r_0\) denote matched no-repair. The
post-repair integrity value at horizon \(\tau\) is
\[
V_t^{\mathrm{rep}}(r;d)
=
\mathbb E
\left[
v_{t+\tau}
\mid
\operatorname{do}(D_t=d),
\operatorname{do}(R_t=r),
\mathbf H_t
\right].
\]
The repair gain is
\[
\Delta_t^{\mathrm{rep}}(r;d)
=
V_t^{\mathrm{rep}}(r;d)
-
V_t^{\mathrm{rep}}(r_0;d).
\]
Repair costs are handled either by restricting the feasible repair family or
by reporting a separate cost-effect curve. They should not be hidden inside
an undocumented scalar reward.

\subtopic{53.7 Damage detection, repair mediation, and specificity}

Let \(M_t\) be an internal integrity monitor and \(R_t\) the selected repair
action. A self-maintenance claim requires the causal chain
\[
D_t
\longrightarrow
M_t
\longrightarrow
R_t
\longrightarrow
v_{t+\tau}.
\]
The required tests are:
\begin{description}
\item[Detection.]
Matched damage interventions change \(M_t\) in a gate-specific direction.

\item[Repair mediation.]
Intervening on \(M_t\) changes the repair distribution, and blocking \(M_t\)
reduces the damage-to-repair effect.

\item[Restoration.]
The selected repair improves \(v_{t+\tau}\) relative to matched no-repair.

\item[Specificity.]
A cost-matched generic action or an off-workspace repair does not produce the
same restoration.

\item[Identity preservation.]
The restored system remains within the declared identity tolerance relative to
the pre-damage workspace.
\end{description}

Define the detection effect \(\eta_{D\to M}\), mediation effect
\(\eta_{M\to R}\), restoration effect \(\eta_{R\to v}\), and repair
specificity
\[
\Delta_{\mathrm{repair}}^{\mathrm{spec}}
=
\eta_{R\to v}
-
\eta_{R_\perp\to v},
\]
where \(R_\perp\) is a matched action outside the proposed repair channel.

\subtopic{53.8 \(T_{21}\): Repair-selection margin theorem }

\begin{theorem}[Repair selection under value-model error]
\label{thm:repair-selection}
Fix a damage intervention \(d\) and a finite feasible repair set
\(\mathcal R\). Let
\[
V(r)=V_t^{\mathrm{rep}}(r;d)
\]
be the true expected post-repair integrity margin and let \(\widehat V\) obey
\[
\sup_{r\in\mathcal R}|\widehat V(r)-V(r)|
\leq
\varepsilon_R.
\]
Let \(\widehat r\in\argmax_r\widehat V(r)\) and
\(r^*\in\argmax_r V(r)\). Then
\[
\boxed{
V(r^*)-V(\widehat r)
\leq
2\varepsilon_R.
}
\]
If
\[
V(r^*)
\geq
\rho+2\varepsilon_R,
\]
then the selected repair restores the expected reserve:
\[
\boxed{V(\widehat r)\geq\rho.}
\]
If \(r^*\) is unique and its repair-value gap exceeds
\(2\varepsilon_R\), then \(\widehat r=r^*\).
\end{theorem}

\begin{proof}
Greedy optimality and uniform approximation give
\[
\begin{aligned}
V(r^*)
&\leq
\widehat V(r^*)+\varepsilon_R\\
&\leq
\widehat V(\widehat r)+\varepsilon_R\\
&\leq
V(\widehat r)+2\varepsilon_R.
\end{aligned}
\]
The reserve and exact-selection statements follow immediately.
\end{proof}

The theorem is conditional on a valid repair-value model. It cannot convert a
post hoc repair narrative into evidence that the system detected damage or
used the model. Detection and value-channel interventions remain necessary.

\subtopic{53.9 Repair, replacement, and successor workspaces}

Restoration of performance is not automatically restoration of identity. Let
\(W_{t^-}\) be the last certified pre-damage workspace and let
\(W_{t+\tau}^{r}\) be the post-repair workspace. Define the repair identity
defect
\[
\delta_{\mathrm{rep-id}}(r)
=
 d_{\mathrm{id}}
\left(
W_{t^-},W_{t+\tau}^{r}
\right).
\]
A repair action is identity preserving at tolerance
\(\varepsilon_{\mathrm{id}}\) when
\[
\Delta_t^{\mathrm{rep}}(r;d)>0,
\qquad
\delta_{\mathrm{rep-id}}(r)
\leq
\varepsilon_{\mathrm{id}}.
\]
When integrity is restored but the identity defect exceeds tolerance, the
result is better described as a successor or replacement workspace. This
terminological distinction prevents any high-performing reconstruction from
being counted as persistence of the same relational organization.

\subtopic{53.10 \(T_{22}\): Closed-loop self-maintenance theorem }

\begin{theorem}[Causal repair stabilizes maintenance deficit]
\label{thm:self-maintenance}
Let \(e_t\geq0\) be the maintenance deficit. Suppose the matched no-repair
trajectory satisfies
\[
e_{t+1}^{(0)}
\leq
\alpha e_t+d_t,
\]
where \(d_t\geq0\) is exogenous damage. Suppose the implemented repair policy
has the causal one-step effect
\[
e_{t+1}^{(0)}-e_{t+1}
\geq
\kappa e_t-\epsilon_t.
\]
Assume
\[
0\leq\beta:=\alpha-\kappa<1,
\qquad
0\leq d_t\leq\overline d,
\qquad
0\leq\epsilon_t\leq\overline\epsilon.
\]
Then
\[
\boxed{
e_t
\leq
\beta^te_0
+
\frac{(\overline d+\overline\epsilon)(1-\beta^t)}{1-\beta},
}
\]
and therefore
\[
\boxed{
\limsup_{t\to\infty}e_t
\leq
\frac{\overline d+\overline\epsilon}{1-\beta}.
}
\]
If the admissible deficit is \(e_{\max}\) and
\[
e_0\leq e_{\max},
\qquad
\overline d+\overline\epsilon
\leq
(1-\beta)e_{\max},
\]
then the interval \([0,e_{\max}]\) is forward invariant.
\end{theorem}

\begin{proof}
Combine the no-repair bound with the causal repair effect:
\[
\begin{aligned}
e_{t+1}
&\leq
 e_{t+1}^{(0)}-\kappa e_t+\epsilon_t\\
&\leq
(\alpha-\kappa)e_t+d_t+\epsilon_t\\
&\leq
\beta e_t+\overline d+\overline\epsilon.
\end{aligned}
\]
Iterating this scalar recurrence gives the stated geometric bound. If
\(e_t\leq e_{\max}\), the forward-invariance condition gives
\[
e_{t+1}
\leq
\beta e_{\max}
+(1-\beta)e_{\max}
=e_{\max}.
\]
\end{proof}

Theorem~\ref{thm:self-maintenance} differs from the persistence theorem in
Part~X. There, contraction was assumed for the workspace error dynamics. Here
the improved contraction factor is attributed to a measured repair action
through a matched no-repair counterfactual. Without that causal comparison,
the same bounded trajectory could be passive robustness rather than
self-maintenance.

\subtopic{53.11 A scalar boundary-repair example}

Let \(b_t\in[0,1]\) be a boundary-integrity score and let the reserve be
\(\rho\). Define
\[
e_t=(\rho-b_t)_+.
\]
Suppose uncorrected damage obeys
\[
e_{t+1}^{(0)}
\leq
1.08e_t+d_t,
\]
so the open-loop boundary is locally unstable. A repair policy with measured
causal effect
\[
e_{t+1}^{(0)}-e_{t+1}
\geq
0.25e_t-\epsilon_t
\]
gives
\[
\beta=1.08-0.25=0.83.
\]
If
\[
\overline d=0.01,
\qquad
\overline\epsilon=0.005,
\]
then
\[
\limsup e_t
\leq
\frac{0.015}{1-0.83}
\approx0.0883.
\]
Thus an open-loop unstable boundary can be maintained within a finite deficit
band by a sufficiently strong repair channel. The example establishes only a
control calculation. A complete self-maintenance claim must also show that an
internal monitor detects damage, mediates the repair choice, and preserves
workspace identity.

\subtopic{53.12 Complete self-maintenance certificate}

A self-maintaining workspace must pass all of the following gates:
\begin{enumerate}
\item \emph{Protected-gate declaration.} The maintained integrity vector and
 reserve are fixed before evaluation.
\item \emph{Damage validity.} Each damage intervention measurably degrades at
 least one protected gate.
\item \emph{Internal detection.} The monitor changes under damage and is not
 merely an external evaluator.
\item \emph{Repair mediation.} Monitor interventions change repair selection.
\item \emph{Repair efficacy.} Repair improves the held-out integrity margin
 relative to matched no-repair.
\item \emph{Specificity.} Cost-matched generic actions and off-workspace
 interventions do not reproduce the gain.
\item \emph{Repeated-disturbance stability.} The deficit remains within the
 declared bound over a held-out disturbance sequence.
\item \emph{Identity control.} The restored state remains within the declared
 identity tolerance or is reported as a successor workspace.
\item \emph{Anti-centralization.} No single node contains a sufficient copy
 of the monitor, repair model, and protected state when the claim is
 workspace-level maintenance.
\end{enumerate}

\subtopic{53.13 Certified workspace agent}

Let
\[
\Theta_A
=
\left(
\tau_G,
\tau_{G\to A},
\varepsilon_Q,
\tau_{V\to A},
\tau_{A\to Y},
\tau_{D\to M},
\tau_{M\to R},
\tau_{R\to v},
\tau_{\mathrm{spec}},
\varepsilon_{\mathrm{ctl}},
 e_{\max}
\right).
\]
A \(\Theta_A\)-certified workspace agent is a controlled realization for
which:
\begin{enumerate}
\item \(\SRW(W_t;\varsigma)=1\);
\item cue-controlled goal retention exceeds \(\tau_G\);
\item goal-to-action mediation exceeds \(\tau_{G\to A}\);
\item counterfactual-value error is at most \(\varepsilon_Q\);
\item value-to-action use exceeds \(\tau_{V\to A}\);
\item action-to-outcome control exceeds \(\tau_{A\to Y}\);
\item damage detection, monitor-to-repair mediation, and restoration exceed
 their thresholds;
\item repair specificity exceeds \(\tau_{\mathrm{spec}}\);
\item \(k_{\mathrm{ctl},\varepsilon_{\mathrm{ctl}}}\geq2\); and
\item repeated maintenance deficit remains at most \(e_{\max}\).
\end{enumerate}

As throughout the theory, this predicate is relative to the declared node
decomposition, goals, actions, environment, horizons, interventions, and null
families.

\subtopic{53.14 \(T_{23}\): Workspace-agency synthesis theorem }

\begin{theorem}[Workspace-level agency under declared gates]
\label{thm:agency-synthesis}
Assume a controlled realization satisfies:
\begin{enumerate}
\item the operational relationality assumptions of
 Theorem~\ref{thm:operational-relationality};
\item cue-controlled goal retention and action readout as in
 Theorem~\ref{thm:goal-retention} and
 Corollary~\ref{cor:goal-action-readout};
\item a counterfactual evaluator with uniform error \(\varepsilon_Q\) and
 optimization error \(\delta_{\mathrm{opt}}\);
\item feasible interventions establishing positive value-to-action and
 action-to-outcome effects;
\item the causal self-maintenance assumptions of
 Theorem~\ref{thm:self-maintenance}; and
\item \(k_{\mathrm{ctl},\varepsilon}\geq2\), together with a matched local
 policy-null panel that preserves every single-node predictive marginal
 but loses full-system control value.
\end{enumerate}
Then, relative to the declared specification:
\begin{enumerate}
\item the expected one-step decision regret is at most
 \(2\varepsilon_Q+\delta_{\mathrm{opt}}\);
\item the long-run maintenance deficit is bounded by
 \[
 \frac{\overline d+\overline\epsilon}{1-\beta};
 \]
\item no single-node statistic is sufficient for the complete control state;
\item matched local policy nulls cannot reproduce the intact control value;
\item goal-dependent action selection and workspace repair are causally
 mediated by implemented internal channels.
\end{enumerate}
Hence the realization satisfies the formal predicate
\[
\operatorname{Agent}_{\Theta_A}(\Workspace)=1.
\]
\end{theorem}

\begin{proof}
The regret bound follows from
Theorem~\ref{thm:counterfactual-policy}. The maintenance bound follows from
Theorem~\ref{thm:self-maintenance}. The coalition condition excludes
single-node sufficiency for the complete control state. The matched-null gap
excludes explanations based only on preserved local predictive marginals. The
implemented goal, value, action, monitor, and repair interventions establish
the required causal mediation. These are precisely the gates in the stated
agency predicate.
\end{proof}

The synthesis theorem is intentionally relative and conjunctive. Its force is
that decision quality, maintenance stability, causal mediation, and relational
location are certified together. It does not infer an unrestricted agent
outside the tested environment.

\subtopic{53.15 Agency, autonomy, and consciousness}

The following implications are not licensed:
\[
\operatorname{Agent}_{\Theta_A}(\Workspace)
\centernot\Longrightarrow
\operatorname{PhenomenallyConscious}(\Workspace),
\]
\[
\operatorname{SelfMaintaining}(\Workspace)
\centernot\Longrightarrow
\operatorname{SelfAware}(\Workspace),
\]
\[
\operatorname{GoalPersistent}(\Workspace)
\centernot\Longrightarrow
\operatorname{MoralPatient}(\Workspace).
\]
A thermostat-like controller can maintain a variable without representing
itself. A sophisticated planner can evaluate counterfactual actions without
phenomenal experience. Conversely, failure to pass this engineering
certificate would not prove absence of consciousness. Agency and
consciousness remain distinct axes of inquiry.
